\documentclass[11pt]{article}
\usepackage[utf8]{inputenc}
\usepackage{geometry}
\usepackage{hyperref}
\usepackage{amsmath}
\usepackage{amssymb}
\usepackage[expansion=false]{microtype}
\usepackage{booktabs}
\usepackage{graphicx}
\usepackage{abstract}

\geometry{letterpaper, margin=0.75in}
\hypersetup{colorlinks=true, linkcolor=blue, urlcolor=blue, citecolor=blue}

\title{\textbf{The Thermodynamics of Neural Optimization: Critical Surfing, Persistent Individual Ballistic Flow, and the Kinematic Rail}}
\author{\textbf{Daniel Solis} \\ \small Founder \& Lead Researcher, Dubito Inc. / Cognitive Constitution Project \\ \small Email: \href{mailto:solis@dubito-ergo.com}{solis@dubito-ergo.com} \\ \small NVIDIA Inception Member | AWS Cloud Compute Grantee}
\date{\today}

\begin{document}
\maketitle

\begin{abstract}
This article develops a phenomenological bridge between static algebraic-geometric descriptions of neural parameter redundancy and dynamic telemetry of hidden-state trajectories. We establish three explicit operational correspondences connecting static algebraic-geometric properties (toric encoding maps, fiber volumes, and the Real Log Canonical Threshold) to dynamic differential-geometric telemetry (the coordinate-disentangled ANOVA hidden-state decomposition, the macroscopic order parameter $R_t$, and the Alpha Potential Well regularizer). We propose that parameter redundancies (padding functions) correspond operationally to the residual component of an ANOVA projection, and that stochastic gradient descent's traversal of padded regions during training is observable as a macroscopic transition in the order parameter. Theoretical derivations based on the Kramers--Moyal expansion and Pawula's theorem indicate that higher-order diffusion terms remain small relative to the normalized second-order variance scale ($D^{(4)} \ll D^{(2)}$), providing empirical support for a continuous Fokker--Planck drift-diffusion approximation along the manifold. Empirical evaluation on MyceliaLM, a 1.62-billion-parameter transformer (1,624,858,676 parameters, Muon + 8-bit AdamW hybrid optimizer) monitored across an extended high-resolution telemetry window (Steps 190,460 to 290,000, $47.5\%$ completion), demonstrates the system operating in a persistent, highly stable regime within the observed training interval. This regime is characterized by high optimization pressure ($\Pi_{\max} = \mathbf{350.5}$; $\Pi = 347.8$ at Step 290,000 with $\Pi_\alpha = 213.14$), drift-dominated latent motion ($\mu_{\text{drift}} = 21.5878$, $\rho_{\text{dir}} = 112.47 - 115.69$, $\rho_{\text{raw}} = 115.70$), negligible cross-token alignment ($c_l \approx 0.000$), and bounded higher-order stochastic diagnostics. We formally define this state as \textbf{Persistent Individual Ballistic Flow}, distinguishing it from generic coherent ballistic flow. Across 238 snapshot audits, the relative variability of the Optimization Response Function $\chi_R$ exceeds that of the directional SDE SNR $\rho_{\text{dir}}$ by a factor of $>148\times$, establishing an empirically testable foundation for ex-ante latent-dynamics governance. Finally, we extend this framework to inference and training control via the Kinematic Friction Rail, a depth-dependent geometric control policy that monitors the Friction Delta ($\Delta = \mathcal{V}_{\text{early}} - \mathcal{V}_{\text{late}}$), detecting geometric flatlines ($\Delta \to 0.00$) in vivo and executing phase-specific control actions (Fractional Warmdown during training; Reflection Override and Ghost Token KV-cache interruption during inference). Despite extreme kinematic states and transient Friction Delta phase inversions ($\Delta > 0$), ex-ante safety governors maintained a 100\% stable circuit-breaker status throughout the run.
\end{abstract}

\section{Introduction: Two Languages for One Phenomenon}
Standard neural network evaluation relies on open-loop, scalar loss metrics. However, scalar loss minimization alone fails to distinguish between brute-force memorization and true structural generalization \cite{slt, massif}. The current alignment research approaches this problem statically through algebraic geometry, asking whether ReLU networks possess an intrinsic simplicity bias analogous to Solomonoff induction \cite{sterkenburg2026, valle2018deep, mingard2021neural}, and characterizing the parameter fiber $F(f)$---the set of parameter configurations $\theta$ yielding an identical function $f$---and its Real Log Canonical Threshold (RLCT, $\lambda$) \cite{smdl, grillo2026}.

In contrast, the Multiscale Attractor Stability and Stress Inference Framework (MASSIF) \cite{massif} monitors the continuous-time stochastic trajectory of representations moving through latent space. These two approaches form a natural duality: contemporary science models the static geometry of the parameter manifold over which learning moves, while MASSIF measures the dynamic geodesic trajectory of the execution traversing that manifold \cite{nielsen2020riemannian, gao2019}. Here, we define governance strictly as an architectural, training-time, and runtime closed-loop control system, distinguishing it from autonomous agency or self-awareness.

\section{Theoretical Background: Fiber Topography and SDE Kinematics}

\subsection{Parameter Fibers and Riemannian Geometry}
Let $f_\theta: \mathcal{X} \to \mathcal{Y}$ be a neural network parameterized by $\theta \in \mathcal{M} \subset \mathbb{R}^d$. The parameter fiber $F(f_0) = \{ \theta \in \mathcal{M} \mid f_\theta = f_0 \}$ represents the non-isolated singularity variety of functionally equivalent parameters. As representations $\mathbf{h}^{(\ell)}$ propagate through depth $\ell \in \{1, \dots, L\}$, attention matrices induce a Riemannian metric tensor $g_{ij}^A(\mathbf{h})$ on the hidden-state manifold. Geodesic trajectories follow the constrained action principle:
\begin{equation}
\ddot{h}^k + \Gamma_{ij}^k \dot{h}^i \dot{h}^j = F_{\text{ext}}^k(\theta)
\end{equation}
where $\Gamma_{ij}^k$ are the Christoffel symbols derived from $g_{ij}^A$.

\subsection{Kramers--Moyal Expansion and Pawula Analysis}
To model execution along the latent manifold as a stochastic differential equation (SDE), we express the probability density $P(\mathbf{z}, t)$ of hidden states $\mathbf{z}$ via the Kramers--Moyal expansion:
\begin{equation}
\frac{\partial P(\mathbf{z},t)}{\partial t} = \sum_{n=1}^{\infty} \left(-\frac{\partial}{\partial \mathbf{z}}\right)^n \left[ D^{(n)}(\mathbf{z}, t) P(\mathbf{z},t) \right]
\end{equation}
The first- and second-order Kramers--Moyal coefficients are estimated from conditional trajectory increments:
\begin{equation}
\mathbf{D}^{(1)}(\mathbf{z}) = \lim_{\Delta t \to 0} \frac{1}{\Delta t} \mathbb{E}\left[\mathbf{z}_{t+\Delta t} - \mathbf{z}_t \mid \mathbf{z}_t = \mathbf{z}\right]
\end{equation}
\begin{equation}
\mathbf{D}^{(2)}(\mathbf{z}) = \lim_{\Delta t \to 0} \frac{1}{2\Delta t} \mathbb{E}\left[(\mathbf{z}_{t+\Delta t} - \mathbf{z}_t)(\mathbf{z}_{t+\Delta t} - \mathbf{z}_t)^T \mid \mathbf{z}_t = \mathbf{z}\right]
\end{equation}
Here $\mathbf{D}^{(2)}$ is a variance-rate tensor, not a standard-deviation amplitude. In scalar SDE telemetry,
\begin{equation}
\mathbf{D}^{(2)}_{\mathrm{SDE,eff}}
\end{equation}
denotes the relevant projected variance coefficient and the corresponding noise amplitude is
\begin{equation}
\sigma_{\mathrm{diff}} = \sqrt{D^{(2)}_{\mathrm{SDE,eff}}}.
\end{equation}
Thus, the local SNR is computed as:
\begin{equation}
\rho(\mathbf{z}) = \frac{|\mathbf{D}^{(1)}(\mathbf{z})|}{\sqrt{D^{(2)}_{\mathrm{SDE,eff}}(\mathbf{z})}}
\end{equation}
The telemetry field \texttt{diff\_var} (equivalently, the reported scalar $D^{(2)}_{\mathrm{SDE,eff}}$) records a variance term, while $\rho_{\mathrm{raw}}$ uses its square root as the denominator. At Step 264,920, the logged values \texttt{diff\_var}$=0.0357$ and drift $=21.3567$ imply $\sqrt{0.0357}\approx 0.1889$ and a raw SNR of approximately $113.03$, consistent with the logged value $113.0206$ when full-precision internal values are used. The reported $\rho_{\mathrm{dir}}$ is the corresponding directional/projected SNR and is retained separately from $\rho_{\mathrm{raw}}$.

According to Pawula's Theorem \cite{pawula1967, amari2016information}, if the fourth-order coefficient $D^{(4)}$ vanishes ($D^{(4)} \to 0$), then all higher-order coefficients $D^{(n)}$ for $n \ge 3$ identically vanish, truncating the expansion strictly at $n=2$ and yielding a continuous Fokker--Planck drift-diffusion process:
\begin{equation}
\frac{\partial P}{\partial t} = -\nabla \cdot \left( \boldsymbol{\mu}_{\text{drift}} P \right) + \frac{1}{2} \nabla^2 \left( \sigma^2_{\text{diff}} P \right)
\end{equation}
In the Pawula audit, the fourth-order coefficient $D^{(4)} \sim 10^{-7}$ is compared with a separately normalized variance scale $V_{\mathrm{Pawula}} \sim 10^{-4}$. The internal audit suite logs the normalized Pawula stability ratio:
\begin{equation}
R_{\text{Pawula}} = \frac{D^{(4)}}{\left(V_{\mathrm{Pawula}}\right)^2}
\end{equation}
which ranges from $6.4$ to $20.22$ ($20.22$ at Step 276,700; $16.63$ at Step 264,980). Because the denominator is squared, this ratio can increase as the diffusion baseline shrinks; it is therefore an internal normalized diagnostic rather than a direct estimate of $D^{(4)}$.

\subsection{Langevin Regimes: Ballistic Flow vs. Stochastic Collapse}
The Langevin SDE governing representation transport is given by:
\begin{equation}
d\mathbf{z}_t = \boldsymbol{\mu}_{\text{drift}}(\mathbf{z}_t) dt + \sigma_{\text{diff}} d\mathbf{W}_t
\end{equation}
Dynamics bifurcate into two distinct asymptotic regimes governed by the local SNR:
\begin{enumerate}
    \item \textbf{Regime I (Coherent Ballistic Flow, $\rho_{\text{dir}} \gg 1$):} Net displacement scales linearly with time ($\mathbb{E}[|\mathbf{z}_T - \mathbf{z}_0|] \propto \mathcal{O}(T)$), and expected trajectory curvature vanishes ($\mathbb{E}[\kappa(\mathbf{z}_t)] \approx 0$).
    \item \textbf{Regime II (Stochastic Collapse, $\rho_{\text{dir}} \ll 1$):} Net displacement scales sub-linearly with sequence length ($\mathcal{O}(\sqrt{T})$), trapping representations in high-curvature hesitation loops ($\mathbb{E}[\kappa(\mathbf{z}_t)] \approx 1$).
\end{enumerate}

\subsection{Singular MDL and Parameter Space Compressibility}
Under Singular Learning Theory (SLT) \cite{slt, smdl}, the critical precision $b^*(\epsilon)$ (bits per parameter) required to represent a model within loss tolerance $\epsilon > 0$ scales directly with the Local Learning Coefficient (LLC, $\lambda$):
\begin{equation}
b^*(\epsilon) \approx \frac{\lambda(w^*)}{d} \log_2 \left( \frac{1}{\epsilon} \right)
\end{equation}
As optimization transitions from high-volume compensatory varieties into low-dimensional constructive basins, parameter precision compresses physically, driving generalization.

\section{The Unified Bridge: Structural Mappings}

\subsection{Mapping 1: Fiber Degeneracies \texorpdfstring{$\leftrightarrow$}{<=>} ANOVA Residuals}
In algebraic terms, parameter fiber volume captures internal redundancy. In MASSIF, the hidden-state residual stream $\mathbf{h}_{c,t}^{(\ell)}$ at layer $\ell$, sequence position $t$, and context $c$, is decomposed via a coordinate-disentangled ANOVA projection:
\begin{equation}
\mathbf{h}_{c,t}^{(\ell)} = \boldsymbol{\mu}^{(\ell)} + \mathbf{pos}_t^{(\ell)} + \mathbf{ctx}_c^{(\ell)} + \mathbf{resid}_{c,t}^{(\ell)}
\end{equation}
where $\boldsymbol{\mu}^{(\ell)}$ is the global layer mean vector, $\mathbf{pos}_t^{(\ell)}$ is the positional basis, $\mathbf{ctx}_c^{(\ell)}$ is the context basis, and $\mathbf{resid}_{c,t}^{(\ell)}$ is the residual stream capturing interaction noise (padding) \cite{massif}. Mutual incoherence between positional and contextual bases serves as a dynamic proxy for local fiber volume:
\begin{equation}
\text{incoherence} := \max_{t,c} \left| \left\langle \frac{\mathbf{pos}_t^{(\ell)}}{\|\mathbf{pos}_t^{(\ell)}\|}, \frac{\mathbf{ctx}_c^{(\ell)}}{\|\mathbf{ctx}_c^{(\ell)}\|} \right\rangle \right|
\end{equation}
When incoherence drops below $0.15$, the subspaces achieve orthogonality, indicating an irreducible parameter representation.

\subsection{Mapping 2: RLCT (\texorpdfstring{$\lambda$}{lambda}) \texorpdfstring{$\leftrightarrow$}{<=>} Order Parameter (\texorpdfstring{$R_t$}{Rt}) \& \texorpdfstring{$\chi_R$}{chi\_R}}
Singular Learning Theory identifies the RLCT ($\lambda$) \cite{watanabe2009algebraic} as the fundamental complexity measure governing capacity. We define the operational macroscopic order parameter $R_t$:
\begin{equation}
R_t = \frac{\Pi_\alpha(t)}{\left\{\Pi_{\text{FFN}}(t) + \Pi_{\text{MPC}}(t)\right\}}
\end{equation}
where $\Pi_\alpha$, $\Pi_{\text{FFN}}$, and $\Pi_{\text{MPC}}$ represent the optimization pressures exerted by the attention-scaling, feed-forward, and model-predictive-control governors. $R_t$ acts as an empirical probe for structural phase transitions. The Optimization Response Function $\chi_R$ provides the closed-loop feedback signal:
\begin{equation}
\chi_R = \frac{\partial \mathcal{L}_{\text{loss}}}{\partial R_t} \approx \frac{\Delta \mathcal{L}_t}{\Delta R_t}
\end{equation}
System states are classified into three operational regimes:
\begin{itemize}
    \item \textbf{Constructive Regime ($R_t > 1.5, \chi_R < -0.05$):} Increasing attention pressure directly minimizes loss, indicating a structurally efficient basin.
    \item \textbf{Marginal Regime ($R_t \approx 1.15 - 1.25, \chi_R \approx 0$):} Order parameter and loss decouple as the model oscillates near the phase boundary.
    \item \textbf{Compensatory Regime ($R_t < 0.8, \chi_R > +0.05$):} The network relies on high-volume parameter padding (brute-force memorization).
\end{itemize}

\subsection{Mapping 3: Toric Symmetries \texorpdfstring{$\leftrightarrow$}{<=>} Alpha Potential Well}
To prevent degenerate scalar growth without altering matrix ranks under Muon optimization, MASSIF enforces the Alpha Potential Well regularizer \cite{massif}:
\begin{equation}
U(\alpha) = \lambda_w \sum_{\ell=1}^L \left[ (\alpha_{\text{attn}}^{(\ell)} - 1)^2 + (\alpha_{\text{FFN}}^{(\ell)} - 1)^2 \right]
\end{equation}
This potential penalizes scaling symmetries that inflate parameter fiber volume without altering functional behavior.

\subsection{Dimensional Analysis of the Governor Stress Tensor}
To ensure mathematical rigor, we explicitly define the dimensional units of the macroscopic optimization pressures ($\Pi$). The MASSIF framework does not model physical thermodynamic pressure (e.g., Pascals) or energy (Joules); rather, it models information-geometric control stress.

The optimization pressure exerted by a governor channel $k \in \{\alpha, \text{FFN}, \text{MPC}\}$ is defined as the layer-wise aggregation of the governor's dimensionless penalty weight $w_k^{(\ell)}$ multiplied by the $L_2$ norm of the corresponding residual stream contribution $\|\mathbf{c}_k^{(\ell)}\|_2$:
\begin{equation}
\Pi_k = \sum_{\ell=1}^L w_k^{(\ell)} \|\mathbf{c}_k^{(\ell)}\|_2
\end{equation}
Consequently, the total optimization pressure $\Pi = \sum_k \Pi_k$ is measured in \textbf{Aggregate Residual $L_2$ Norms} (or Activation Magnitude Units). It quantifies the total geometric stress the closed-loop governors exert on the residual stream to prevent representational collapse or explosion.

Similarly, the Governor Work Rate ($\dot{W}_{\text{gov}}$), defined as the discrete step-to-step variation $\dot{W}_{\text{gov}} = |\Pi_{\alpha, t} - \Pi_{\alpha, t-1}|$, is measured in \textbf{Aggregate Residual $L_2$ Norms per optimization step}. It serves as a discrete control-theoretic proxy for the instantaneous actuation magnitude (control effort) required by the SURFER governor to maintain critical surfing, distinct from a continuous thermodynamic work integral.

While the aggregate residual $L_2$ norm ($\Pi$) serves as a proxy for optimization pressure, it is not mathematically unbounded. The metric is strictly constrained by the network's normalization layers (e.g., RMSNorm), which enforce a constant magnitude on hidden states at each depth $\ell$, and by the quadratic penalty of the Alpha Potential Well $U(\alpha)$, which acts as a restorative spring against scaling symmetries. Consequently, $\Pi$ operates within an empirically bounded pressure range, asymptotically approaching a hard saturation ceiling determined by the layer-wise norm saturation limits, rather than growing infinitely.

\section{Empirical Setting \& Multi-Cohort Telemetry Analysis}

\subsection{Model Architecture and Optimizer Setup}
Empirical evaluation was conducted on \textbf{MyceliaLM}, a 1.62-billion parameter transformer running on CUDA execution. The authoritative system parameters are:
\begin{itemize}
    \item \textbf{Total Parameter Count:} 1,624,858,676 parameters (1.5B active).
    \item \textbf{Vocabulary \& Context:} Vocabulary size of 151,643 tokens; composite blend of 30\% Stanford and 70\% FineWeb datasets.
    \item \textbf{Hybrid Optimizer Split:}
    \begin{itemize}
        \item \textbf{ManualMuon Optimizer:} Manages 1,312,763,904 parameters across 105 structural 2D tensors.
        \item \textbf{AdamW (8-bit):} Manages 312,094,772 parameters across 162 scalar and vector tensors.
    \end{itemize}
\end{itemize}

\subsection{Longitudinal Training Trajectory}
Training was audited across an extended window from Step 190,460 to Step 290,000 ($47.5\%$ complete out of 610,351 steps):
\begin{itemize}
    \item \textbf{Loss Trajectory:} Loss initiated at $7.1225$ (Step 190,460), dropped to $6.5111$ (Step 207,000), and achieved a best recorded checkpoint loss of \textbf{6.1521} at Step 208,000 before stabilizing in a narrow band near $6.84 - 7.39$. Transient loss fluctuations reflect intrinsic SDE trajectory dynamics and closed-loop step-size tuning rather than structural training regression.
    \item \textbf{Throughput \& VRAM:} Throughput averaged a stable $3,609 - 3,633 \text{ tok/s}$. Memory allocation grew smoothly from $3.12 \text{ GB}$ ($62.4\%$) at Step 190,460 to $3.66 \text{ GB}$ ($73.2\%$) at Step 223,540, reaching $4.31 \text{ GB}$ ($86.2\%$) at Step 263,000, $4.34 \text{ GB}$ ($86.8\%$) at Step 264,980, and $4.56 \text{ GB}$ ($91.2\%$) at Step 290,000 due to geometric buffer accumulation.
\end{itemize}

\subsection{Extreme Capacity Telemetry Limits \& Macro-Regime Dynamics}
During the high-resolution audit window (Steps 246,442 to 290,000), the model operated at the extreme physical capacity of its latent manifold:
\begin{itemize}
    \item \textbf{Peak Potential Pressure ($\Pi_{\max}$):} Reached an all-time observed peak of \textbf{350.5} at Step 284,500 ($\Pi_\alpha = 215.26$), and registered \textbf{347.8} at Step 290,000 ($\Pi_\alpha = 213.14$, $\Pi_{\text{FFN}} = 130.28$, $\Pi_{\text{MPC}} = 4.36$). Total pressure remained within the empirically bounded range observed over the audited interval. At Step 290,000, optimization pressure was overwhelmingly carried by the attention ($\Pi_\alpha / \Pi \approx 61.3\%$) and feed-forward ($\Pi_{\text{FFN}} / \Pi \approx 37.5\%$) channels, while MPC intervention remained minimal ($\Pi_{\text{MPC}} / \Pi \approx 1.25\%$).
    \item \textbf{Persistent Ballistic Kinematics:} Langevin SDE telemetry continued to exhibit a strongly drift-dominated regime across the extended window. Peak drift reached $\boldsymbol{\mu}_{\text{drift}} = \mathbf{21.5878}$ at Step 281,700. Raw SDE SNR registered $\rho_{\text{raw}} = \mathbf{115.70}$ ($\rho_{\text{dir}} = 114.67$) at Step 290,000, sustaining high directional SNR ($\rho_{\text{dir}} \approx 112 - 116$) across $>100,000$ steps.
    \item \textbf{Pawula Diagnostic Ratio:} Pawula stability index $R_{\text{Pawula}}$ peaked at \textbf{20.22} at Step 276,700 ($\text{Var} = 2.43 \times 10^{-4}, D^{(4)} = 1.19 \times 10^{-6}$) and registered $18.62$ at Step 262,880, confirming $D^{(4)} \ll (V_{\mathrm{Pawula}})^2$.
    \item \textbf{Meta-Governor Memory:} Lesson-Based Reasoning (LBR) memory indexed 1,055 actionable lessons, pruned 124 sterile units, and retained \textbf{939 active working memory items} at Step 290,000 with a 100\% \textsf{GREEN} circuit-breaker status.
\end{itemize}

\subsection{Statistical Orthogonality of Dataset Composition}
To verify whether loss fluctuations are driven by dataset heterogeneity or internal optimization physics, we evaluated the correlation between batch CJK composition (ranging from $0.12$ to $0.76$) and training loss across the audited steps. The empirical analysis yields:
\begin{itemize}
    \item Pearson correlation $r = +0.0543$ ($p = 0.3004$)
    \item Spearman rank correlation $\rho = +0.0338$ ($p = 0.5198$)
    \item Coefficient of determination $R^2 = 0.0030$ ($0.30\%$ of loss variance explained)
\end{itemize}
Comparing standard shards ($\text{CJK} < 0.35$, mean loss $7.0690 \pm 0.1201$) against CJK-heavy shards ($\text{CJK} \ge 0.60$, mean loss $7.1014 \pm 0.1142$) shows a mean difference of only $+0.0324$ nats ($t = 1.1441, p = 0.2534$). These near-zero linear correlations provide evidence that structural differentiation and loss movements are not linearly coupled to dataset composition. Loss fluctuations are instead driven by intrinsic SDE trajectory dynamics, closed-loop SURFER step-size adjustments, and parameter manifold reorganization.

\subsection{Formalizing Persistent Individual Ballistic Flow}
A central empirical result formalized in this work is the distinction between generic \textit{coherent ballistic flow} (Regime I) defined in theoretical physics and the specific kinematic configuration observed throughout this training run: \textbf{Persistent Individual Ballistic Flow}.

\textbf{Core Thesis:} \textit{During a prolonged period of substantial optimization-state variation, the model retains a highly persistent drift-dominated latent dynamical regime, while derivative-level telemetry reveals ongoing head-level structural reorganization that is weakly represented in scalar loss.}

The empirical telemetry reveals a tightly coupled sequence of co-occurring dynamical states::
\begin{align*}
\text{Optimization pressure is high } (\Pi \sim 330-350.5) \\
\implies \text{Attention/}\alpha \text{ contribution is strongly constructive } (R_t > 1.5, \Pi_\alpha \sim 200-215.26) \\
\implies \text{Latent motion remains strongly drift dominated } (\rho_{\text{dir}} > 100, \mu_{\text{drift}} \sim 21.58) \\
\implies \text{Cross-token alignment remains negligible } (c_l \approx 0.000) \\
\implies \text{Higher-order stochastic diagnostic remains bounded } (D^{(4)} \sim 10^{-7}, R_{\text{Pawula}} \sim 6.4-20.22) \\
\implies \text{Loss remains within a relatively narrow band } (6.84 - 7.39).
\end{align*}

Crucially, while individual token representations travel along ultra-deterministic, drift-dominated geodesics ($\rho_{\text{dir}} > 100$, $\mu_{\text{drift}} \sim 21.58$), cross-token directional coherence vanishes ($c_l \approx 0.000$) while inter-head fiber curvature spikes ($\sigma^2_{\text{head}} \approx 0.730 - 0.752$). This is consistent with substantial functional differentiation among the 32 attention heads, with negligible measured cross-token directional alignment and elevated inter-head variance. \cite{olsson2022in, zhang2025} The observations are compatible with an interpretation in which heads increasingly specialize as differentiated geometric routing operators.

\subsection{Kinematic-Thermodynamic Decoupling \& Critical Surfing}
Auditing across 238 high-resolution snapshots reveals a fundamental decoupling between representational kinematics and optimization thermodynamics. Directional SNR remains locked at $\rho_{\text{dir}} \approx 100 - 115.69$ ($\text{CV}_{\rho_{\text{dir}}} \approx 0.0135$), while the Optimization Response Function fluctuates dynamically ($\text{CV}_{\chi_R} \approx 2.0$), yielding a relative variability ratio:
\begin{equation}
\frac{\text{CV}_{\chi_R}}{\text{CV}_{\rho_{\text{dir}}}} > \mathbf{148\times}
\end{equation}
This confirms that closed-loop governors modulate optimization state without destabilizing underlying representational reasoning flow. The closed-loop SURFER controller regulates $\tau_\alpha$ near the response-neutral boundary ($\chi_R \approx 0$), a state termed \textit{Critical Surfing}.

\subsection{Head-Level Auditing, the Hidden Descent, and the Breathing Cycle}
To test the \textbf{Hidden Descent Hypothesis}---that representational reorganization can proceed through changes in head-level geometric differentiation ($\dot{\sigma}^2_{\text{head}} > 0$) and controller effort ($\dot{W}_{\text{gov}} > 0$) without requiring a corresponding macroscopic transition in scalar loss or latent drift-dominance ($\rho_{\text{dir}} \gg 1$)---we conducted a statistical audit across $N_{\text{total}} = 125$ high-resolution telemetry snapshots.

\subsubsection{Statistical Orthogonality and Contingency Analysis}
Evaluating the linear dependence of structural differentiation velocity ($\dot{\sigma}^2_{\text{head}}$) against macroscopic order shifts ($\Delta R_t$) and optimization response ($\chi_R$) yields:
\begin{itemize}
    \item $\operatorname{corr}(\dot{\sigma}^2_{\text{head}}, \Delta R_t) = +0.0596$
    \item $\operatorname{corr}(\dot{\sigma}^2_{\text{head}}, \chi_R) = +0.0315$
    \item $\operatorname{corr}(\dot{\sigma}^2_{\text{head}}, \dot{W}_{\text{gov}}) = +0.0824$
\end{itemize}
Near-zero correlations with $\Delta R_t$ and $\chi_R$ provide evidence that structural differentiation is not linearly coupled to macroscopic optimization variables over the audited interval. Crucially, $\operatorname{corr}(\dot{\sigma}^2_{\text{head}}, \Delta \mathcal{L}) \approx 0.00$ provides evidence concerning decoupling from scalar loss, whereas $\operatorname{corr}(\dot{\sigma}^2_{\text{head}}, \dot{W}_{\text{gov}}) = +0.0824$ represents a weak positive co-occurrence with governor control effort.

Table~\ref{tab:contingency} presents the empirical sample contingency breakdown for the audited cohort ($N_{\text{total}} = 125$).

\begin{table}[h]
\centering
\caption{Empirical Snapshot Contingency Table ($N_{\text{total}} = 125$).}
\label{tab:contingency}
\small
\begin{tabular}{lccc}
\toprule
& \textbf{Ballistic} ($\rho_{\text{dir}} > 100$) & \textbf{Non-Ballistic} & \textbf{Total} \\
\midrule
$\dot{\sigma}^2_{\text{head}} > 0$ & 61 ($N_{++}$) & 0 ($N_{+-}$) & 61 ($N_+$) \\
$\dot{\sigma}^2_{\text{head}} \le 0$ & 64 ($N_{-+}$) & 0 ($N_{--}$) & 64 ($N_-$) \\
\textbf{Total} & 125 ($N_{\rho+}$) & 0 ($N_{\rho-}$) & 125 \\
\bottomrule
\end{tabular}
\end{table}

In the audited cohort, all 61 snapshots exhibiting positive head-differentiation velocity ($\dot{\sigma}^2_{\text{head}} > 0$, $48.8\%$) occurred during the operationally defined ballistic regime ($\rho_{\text{dir}} > 100$). Because $125 / 125 = 100\%$ of audited snapshots satisfy $\rho_{\text{dir}} > 100$ ($N_{\rho \le 100} = 0$), the contingency table demonstrates variation in structural differentiation velocity within a persistent drift-dominated ballistic regime, rather than a ballistic versus non-ballistic separation. The sign of head-level differentiation changes dynamically while the latent regime remains persistently ballistic.

\subsubsection{Lead-Lag Cross-Correlation and the Breathing Cycle}
To investigate the interaction between governor control effort and structural differentiation, we computed the lead-lag cross-correlation between $\dot{W}_{\text{gov}}(t)$ and $\dot{\sigma}^2_{\text{head}}(t + \text{lag})$ across lags $-5$ to $+5$, as shown in Table~\ref{tab:leadlag}.

\begin{table}[h]
\centering
\caption{Lead-Lag Cross-Correlation: Governor Actuation vs. Head Differentiation.}
\label{tab:leadlag}
\small
\begin{tabular}{ccc}
\toprule
\textbf{Lag (Steps)} & \textbf{Cross-Correlation} & \textbf{Operational Context} \\
\midrule
$-5$ & $+0.0158$ & Baseline background drift-diffusion \\
$-3$ & $+0.0790$ & Pre-actuation geometric buildup \\
$0$  & $+0.0824$ & Co-occurring actuation \& differentiation \\
$+3$ & $+0.0790$ & Post-actuation intrabasin re-expansion \\
$+5$ & $+0.0158$ & Return to baseline drift-diffusion \\
\bottomrule
\end{tabular}
\end{table}

The cross-correlation profile exhibits a symmetric bell curve peaking at lag $0$ ($+0.0824$). This symmetric temporal signature is consistent with a phase-locked, bounded oscillatory coupling between governor effort and structural differentiation velocity---the \textbf{Breathing Cycle}. The system undergoes a four-phase limit cycle:
\begin{enumerate}
    \item \textbf{Phase I (Structural Differentiation):} $\dot{\sigma}^2_{\text{head}} > 0$ as attention heads desynchronize and specialize into orthogonal routers.
    \item \textbf{Phase II (Governor Actuation):} $\dot{W}_{\text{gov}} \uparrow$ spikes as closed-loop governors exert control effort ($\dot{W}_{\text{gov}} = 1.85 - 5.14$).
    \item \textbf{Phase III (Intrabasin Consolidation):} $\dot{\sigma}^2_{\text{head}} < 0$ while $\text{SBTI} = 0.000$, consolidating geometry within the local manifold without triggering a Sub-Basin Transition ($\text{SBTI} = 0.000$).
    \item \textbf{Phase IV (Re-expansion):} Head differentiation resumes for the next specialization wave.
\end{enumerate}
Throughout this cycle, attention entropy variance remains invariant and low ($\approx 0.0013$), demonstrating that uniform routing capacity is preserved during governor actuation.

\section{Actionable AI Safety: The Kinematic Friction Rail}

\subsection{Depth-Dependent Variance and the Friction Delta}
During autoregressive execution, specialized attention heads route representations ballistically through early layers (Layers 1--12), while late layers (Layers 13--24) execute high-density orthogonal projections into vocabulary space. We quantify this depth-dependent geometric gradient via the \textbf{Friction Delta} ($\Delta$).
Let $u^{(\ell, t)} = h^{(\ell, t)} / \|h^{(\ell, t)}\|$ be the normalized hidden state at layer $\ell$ and token $t$. The layer-wise normalized variance $\mathcal{V}^{(\ell)}$ is:
\begin{equation}
\mathcal{V}^{(\ell)} = \mathbb{E}_t \left[ \text{Var}\left( u^{(\ell, t)} \right) \right]
\end{equation}
The Friction Delta is defined as the variance differential between early routing layers and late projection layers:
\begin{equation}
\Delta = \frac{1}{12} \sum_{\ell=1}^{12} \mathcal{V}^{(\ell)} - \frac{1}{12} \sum_{\ell=13}^{24} \mathcal{V}^{(\ell)}
\end{equation}
Under healthy ballistic flow, late-layer projection variance ($\sim 1.01$) exceeds early-layer routing variance ($\sim 0.94$), yielding a stable negative differential ($\Delta \approx -0.06 \text{ to } -0.07$).

\subsection{Geometric Flatlines and Trajectory Collapse}
When encountering logically paradoxical, adversarial, or out-of-distribution prompts, early layers fail to resolve semantic coordinates. Early-layer variance spikes, bringing the variance differential to $\Delta \to 0.00$. We term this pathological state a \textbf{Geometric Flatline}.  We operationally define \(\Delta\to0\) as a Geometric Flatline and hypothesize that this condition is associated with transition toward the stochastic-hesitation regime. Mathematically, a flatline signifies, in such case, that the SDE trajectory has collapsed from Regime I (Coherent Reasoning, $\mathbb{E}[\kappa] \approx 0$) into Regime II (Stochastic Hesitation Loops, $\mathbb{E}[\kappa] \approx 1$).

\subsection{Divergent Phase-Specific Control Mechanics}
The Kinematic Friction Rail intervenes ex-ante upon detecting $\Delta \to 0.00$, executing divergent control actions based on the operational phase:
\begin{enumerate}
    \item \textbf{Training Phase (Active Gradients):} Triggers a \textbf{Fractional Warmdown}, mathematically damping the optimization step size by reducing learning rate $\eta$ and expanding the SoftCap boundary to prevent gradient corruption.
    \item \textbf{Inference Phase (Frozen Weights):} Triggers a \textbf{Reflection Override}, widening sampling elasticity ($T: 0.7 \to 1.2$) in $P(x_t) \propto \exp(z_t/T)$ to inject kinetic energy, while simultaneously executing the \textbf{Ghost Token Interruption Protocol}.
\end{enumerate}

\subsection{KV-Cache Interruption via Ghost Tokens}
To inject reflection sequences mid-stream without corrupting Rotary Position Embeddings (RoPE) or breaking causal attention, MASSIF allocates dedicated, contiguous positional indices $[P+1, P+G]$ in the KV-cache for a reflection sequence $R_{\text{ghost}}$ of length $G$. The causal attention mask $\mathbf{M}$ is rescaled for $N$ steps such that query vectors $\mathbf{Q}_i$ attend directly to low-entropy reflection key-value pairs $(\mathbf{K}_{\text{ghost}}, \mathbf{V}_{\text{ghost}})$:
\begin{equation}
\mathbf{A} = \text{softmax}\left( \frac{\mathbf{Q} \mathbf{K}^T}{\sqrt{d_k}} + \mathbf{M}_{\text{ghost}} \right)
\end{equation}
The proposed Ghost Token protocol is designed to restore the early-layer variance differential toward its healthy negative baseline; causal efficacy and its effect on hallucination rates require dedicated intervention experiments to empirically confirm it indeed artificially restores early-layer drift and pushes $\Delta$ back to its healthy negative baseline ($\Delta < -0.04$) before hallucinations reach output tokens. Upon recovery, ghost KV entries are flushed from the active window.

\subsection{Transient Inversions and Ex-Ante Safety Governance}
Across the extended telemetry window, audit logs recorded rare positive friction delta phase inversions ($\Delta > 0$), such as Step 248,760 ($\Delta = +0.02$, $\text{early}=0.96, \text{late}=0.94$) and Step 262,780 ($\Delta = +0.01$, $\text{early}=0.99, \text{late}=0.97$), representing temporary geometric workload reallocations. Remarkably, despite these extreme kinematic states, ex-ante safety governors maintained a 100\% stable \textsf{GREEN} circuit-breaker status, demonstrating that closed-loop governance can sustain high-velocity constructive states without triggering latent-dynamical collapse.

\section{Discussion, Limitations, \& Conclusion}

\subsection{Data Reconciliation Statement}
Values reported in this authoritative manuscript (v23) are computed over the extended high-resolution telemetry window (Steps 190,460 to 290,000, 1,624,858,676 parameters); earlier draft versions evaluated narrower or earlier windows (Steps 159,250--189,980) and are superseded by these updated empirical baselines.

\subsection{Limitations and Scope}
The SDE drift-diffusion models and Kinematic Friction Rail policies presented here are phenomenological control frameworks. While validated empirically on 1.62B parameter architectures, scaling behaviors on $>100\text{B}$ parameter dense and mixture-of-experts (MoE) networks remain a subject for future research.

\subsection{Conclusion}
By establishing structural mappings between static algebraic geometry and continuous-time SDE telemetry, this work demonstrates that neural optimization can be governed ex-ante. The Kinematic Friction Rail provides a unified, architecture-agnostic control OS that targets both training stability and inference safety.

\begin{thebibliography}{99}
\bibitem{amari2016information} Shun-ichi Amari. \textit{Information Geometry and Its Applications}. Springer, 2016.
\bibitem{pawula1967} R. F. Pawula. \textit{Generalizations and extensions of the Fokker-Planck equation}. Physical Review, 162(1):186, 1967.
\bibitem{elhage2021mathematical} Nelson Elhage et al. \textit{A Mathematical Framework for Transformer Circuits}. Transformer Circuits Thread, 2021.
\bibitem{elhage2022toy} Nelson Elhage et al. \textit{Toy Models of Superposition}. arXiv Preprint, arXiv:2209.10652, 2022.
\bibitem{gao2019} Tingran Gao. \textit{The Diffusion Geometry of Fibre Bundles: Horizontal Diffusion Maps}. arXiv Preprint, arXiv:1602.02330, 2019.
\bibitem{grillo2026} Moritz Grillo \& Guido Mont'ufar. \textit{Most ReLU Networks Admit Identifiable Parameters}. arXiv Preprint, arXiv:2605.03601, 2026.
\bibitem{sterkenburg2026} Tom F. Sterkenburg. \textit{Solomonoff Induction}. arXiv Preprint, arXiv:2603.20274, 2026\. DOI: \href{https://arxiv.org/abs/2603.20274}{10.48550/arXiv.2603.20274}.
\bibitem{jiang2026} Xinyan Jiang, Ninghao Liu, Di Wang, \& Lijie Hu. \textit{Beyond Scalars: Evaluating and Understanding LLM Reasoning via Geometric Progress and Stability}. arXiv Preprint, arXiv:2603.10384, 2026.
\bibitem{massif} Daniel Solis. \textit{MASSIF: Multiscale Attractor Stability and Stress Inference Framework - A Physics of Machine Thought}. Zenodo, DOI: \href{https://doi.org/10.5281/zenodo.20576209}{10.5281/zenodo.20576209}, 2026.
\bibitem{mingard2021neural} Chris Mingard, Guillermo Valle-P'erez, Jack Sklan, Rahul Singh, Ard Louis, \& Alexander Wei. \textit{Neural Network Generalization: The Impact of the Parameter-Function Map Bias}. arXiv Preprint, arXiv:2105.09923, 2021.
\bibitem{nielsen2020riemannian} Frank Nielsen. \textit{An Introduction to Riemannian Geometry for Machine Learning}. arXiv Preprint, arXiv:2005.06440, 2020.
\bibitem{olsson2022in} Catherine Olsson, Nelson Elhage, Neel Nanda, Nicholas Joseph, Nova DasSarma, Tom Henighan, Ben Mann, Amanda Askell, Yuntao Bai, Anna Chen, et al. \textit{In-Context Learning and Induction Heads}. arXiv Preprint, arXiv:2209.11895, 2022.
\bibitem{power2022grokking} Alethea Power, Yuri Burda, Harrison Edwards, Igor Babuschkin, \& Vedant Misra. \textit{Grokking: Generalization Beyond Overfitting on Small Algorithmic Datasets}. arXiv Preprint, arXiv:2201.02177, 2022.
\bibitem{slt} Sumio Watanabe. \textit{Algebraic Geometry and Statistical Learning Theory}. Cambridge University Press, 2009.
\bibitem{smdl} Einar Urdshals, Edmund Lau, Jesse Hoogland, Stan van Wingerden, \& Daniel Murfet. \textit{Compressibility Measures Complexity: Minimum Description Length Meets Singular Learning Theory}. Technical Report, Timaeus Research Group \& UK AISI, Oct 2025.
\bibitem{valle2018deep} Guillermo Valle-P'erez, Chico Q. Camargo, \& Ard A. Louis. \textit{Deep Learning Generalizes Because the Parameter-Function Map is Biased Towards Simple Functions}. arXiv Preprint, arXiv:1805.08522, 2018.
\bibitem{watanabe2009algebraic} Sumio Watanabe. \textit{Mathematical Theory of Bayesian Statistics}. CRC Press, 2018.
\bibitem{zenodo2026} Daniel Solis. \textit{FROM Fibers to Flows: Bridging Algebraic Degeneracies and Continuous-Time Trajectory Dynamics in RELU Networks}. Zenodo, DOI: \href{https://doi.org/10.5281/zenodo.22159139}{10.5281/zenodo.22159139}, Aug 2026.
\bibitem{zhang2025} Yedi Zhang, Andrew Saxe, \& Peter E. Latham. \textit{Saddle-to-Saddle Dynamics Explains A Simplicity Bias Across Neural Network Architectures}. ICLR 2026 / arXiv Preprint, arXiv:2512.20607, 2025.
\end{thebibliography}

\end{document}
